\section{Exact hulls for chains with few observed labels}\label{sec:fixed-chains}

We return to the flat chain $G_L$ of \cref{sec:sp-growth} with arbitrary
observations on the gadget arcs and the bypass. The common profile $w$
gives a small exact description (\cref{prop:chain-profile}); eliminating its
residual coordinate sharpens the determinant bound and yields unit flow and
product coefficients when $m\le3$.

\subsection{The complete shared-profile system}

Every description below includes the checks $x\in\Flow$,
$y\in\Delta_m$, and $z\ge0$ on the observed coordinates. Put $t=1-x_h$, and index all $d=m+1$ states by $\states$.
For each gadget $i$, partition the states into $A_i,B_i,T_i,U_i$:
only $a_i$ is observed, only $b_i$ is observed, both are observed,
or neither is observed, respectively. In particular $0\in U_i$.
Write $u_{ij}=z_{a_i,j}$ and $v_{ij}=z_{b_i,j}$ wherever present,
and define
\begin{equation}\label{eq:chain-R}
 R_i=x_{a_i}-\sum_{j\in A_i\cup T_i}u_{ij}
                  +\sum_{j\in B_i}v_{ij}.
\end{equation}
The state class $T_i$ is unrelated to the table polytope $\mathcal T$
of \cref{sec:universality}.

\begin{proposition}[Exact profile formulation]\label{prop:chain-profile}
Subject to the preceding original-domain checks, $(x,y,z)\in\Hull$
if and only if there is $w\in\R^{m+1}$ satisfying
\begin{align}
 0\le w_j\le\lambda_j\quad(j\in\states),\qquad
 &\sum_{j\in\states}w_j=t,\label{eq:chain-profile-bounds}\\
 w_j&=\lambda_j-z_{h,j}&&((h,j)\in\Obs),\label{eq:chain-bypass}\\
 w_j&\ge u_{ij}&&(j\in A_i),\nonumber\\
 w_j&\ge v_{ij}&&(j\in B_i),\nonumber\\
 w_j&=u_{ij}+v_{ij}&&(j\in T_i),\label{eq:chain-cells}\\
 \sum_{j\in B_i}w_j\le R_i&\le
       \sum_{j\in B_i\cup U_i}w_j&&(i\in[L]).\label{eq:chain-endpoints}
\end{align}
\end{proposition}
\begin{proof}
Every state disaggregation supplies its profile through
\eqref{eq:chain-profile}. Its observed entries give
\eqref{eq:chain-bypass}--\eqref{eq:chain-cells}. In gadget $i$,
the $a_i$ entries are $u_{ij}$ in $A_i\cup T_i$, are $w_j-v_{ij}$
in $B_i$, and are free in $[0,w_j]$ in $U_i$. Thus their possible
aggregate is exactly the interval
\[
 \left[\sum_{A_i\cup T_i}u_{ij}+\sum_{B_i}(w_j-v_{ij}),\quad
 \sum_{A_i\cup T_i}u_{ij}+\sum_{B_i}(w_j-v_{ij})+
 \sum_{U_i}w_j\right].
\]
Membership of $x_{a_i}$ in this interval is
\eqref{eq:chain-endpoints}.
Conversely these inequalities allow the free $U_i$ entries to be
chosen with the required sum, independently for each gadget. All
entries lie in $[0,w_j]$: for a doubly observed state this follows
from $w_j=u_{ij}+v_{ij}$ and the original nonnegativity check.
Set every $b_i$ entry to $w_j-f^j_{a_i}$ and the bypass entry to
$\lambda_j-w_j$. State capacities and balances hold, and
$x_{a_i}+x_{b_i}=t$ supplies each $b_i$ aggregate. The bypass
aggregate is $\sum_j\lambda_j-t=x_h$. This is a complete
disaggregation. If a weight is zero, its profile and all its entries
are zero, so no positive-weight assumption has been used.
\end{proof}

Eliminate the unobserved residual coordinate by
$w_0=t-\sum_{j=1}^mw_j$, and write $\widehat w=(w_1,\ldots,w_m)$.
The residual bounds become
\begin{equation}\label{eq:chain-residual-reduced}
 \one^\top\widehat w\le t,\qquad
 -\one^\top\widehat w\le\lambda_0-t.
\end{equation}
Since $0\in U_i$, the endpoint rows become
\begin{equation}\label{eq:chain-endpoints-reduced}
 \one_{B_i}^\top\widehat w\le R_i,\qquad
 \one_{A_i\cup T_i}^\top\widehat w\le t-R_i.
\end{equation}
All remaining rows have singleton normals $\pm e_j$.
Thus each nonzero normal belongs to the fixed universe
\begin{equation}\label{eq:chain-normal-universe}
 \mathcal N_m=\{\one_S:\varnothing\ne S\subseteq[m]\}
                 \cup\{-e_j:j\in[m]\}\cup\{-\one\}.
\end{equation}
Repeated normals represent different affine right-hand sides; keep
the minimum of these right-hand sides at a query point. A zero
normal gives a scalar inequality and must be checked separately.
Every original flow or product has coefficient in $\{-1,0,1\}$
in every individual right-hand side, including $t-R_i$.

\subsection{Finite circuits and constructive recovery}

Let $\Delta^{01}_m$ be the maximum absolute determinant of a square
$0/1$ matrix of order at most $m$, including the empty determinant
one. In particular $\Delta^{01}_1=\Delta^{01}_2=1$, $\Delta^{01}_3=2$, and
$\Delta^{01}_m\le\max(1,m^{m/2})$ by Hadamard's inequality.

\begin{theorem}[Fixed-state chain hull and recovery]\label{thm:fixed-chain}
For $m\ge1$, the sparse hull on $G_L$ has a finite
original-coordinate description with integer flow and product
coefficients of magnitude at most $(m+1)\Delta^{01}_m$.
After $2^{O(m^2)}$ preprocessing depending only on $m$, exact
separation and state-flow recovery each require
\[
 O(mL+|\Obs|+m)+2^{O(m^2)}
\]
rational arithmetic operations. Recovery calls no online linear
program. Rational encoding lengths remain polynomial in the input
length. For $m=0$, the hull is simply the original flow polytope.
\end{theorem}
\begin{proof}
Form the reduced profile inequalities
$a^\top\widehat w\le\gamma_a$ for the distinct present nonzero
normals, where $\gamma_a$ is the minimum of all right-hand sides
with normal $a$. There are at most $2^m+m$ such normals.
The system is feasible exactly when all its nonnegative cancellation
certificates have nonnegative right-hand side, by Farkas' lemma.
The cone of such certificates is
\[
 \{\mu\ge0:\ \sum_a\mu_a a=0\}.
\]
It is pointed. Each extreme ray has minimally dependent support,
of size at most $m+1$; otherwise a linear dependence within its
support permits a small perturbation in both directions and contradicts
extremality. Conversely every minimal positive dependence gives such
an extreme ray. We call these supports, with their primitive positive
integer weights, \emph{positive circuits}.

For a support of size $s$, select $s-1$ independent columns in
its row matrix. The signed $(s-1)$-row minors give a nonzero kernel
vector and hence its circuit weights after a common sign change and
division by the greatest common divisor. Each row is a $0/1$ vector
or its negative. Changing whole-row signs bounds every such minor
by $\Delta^{01}_m$. Therefore every primitive weight is at most
$\Delta^{01}_m$. The extreme rays generate the certificate cone, so
the finitely many tests
\begin{equation}\label{eq:chain-circuit-tests}
 \sum_{a\in C}\mu_a\gamma_a\ge0
 \quad\text{for every positive circuit $C$ of present normals}
\end{equation}
are necessary and sufficient, together with the scalar zero-row
checks. Each circuit is one in the full fixed universe
\eqref{eq:chain-normal-universe}.

Expanding \eqref{eq:chain-circuit-tests} over all choices of an
original row for each participating normal gives a finite affine
description. Positivity of the weights makes this expansion exact:
the weighted sum of minima is the minimum over independent choices.
Each affine inequality sums at most $m+1$ rows, with weights at
most $\Delta^{01}_m$, and hence has the claimed original-coordinate
coefficient bound. At an infeasible query it suffices to output the
rows attaining the minima in one violated circuit. This produces
a globally valid affine inequality, even though other rows may
attain the minima at other points. The original-domain prechecks
have unit coefficients as well.

Enumerate all supports of size at most $m+1$ in the fixed universe,
and use exact elimination to keep their positive circuits. There
are $2^{O(m^2)}$ candidates. Constructing and grouping the profile
rows takes $O(mL+|\Obs|+m)$ arithmetic operations; the circuit
tests then have parameter-only cost $2^{O(m^2)}$.
This count assumes that subset normals are preindexed by bit masks and
singleton indices are cached. A scan of each gadget builds its two
subset masks and right-hand sides, and each local cell row uses a
cached index, so no $m$-entry normal is copied for a local row. The
masks have $O(m)$ bits, which the bit-length bound below includes.

For recovery, also precompute inverses of all nonsingular matrices
formed by $m$ distinct normals of $\mathcal N_m$. At a feasible
query enumerate those whose rows are present and solve their
equalities against the corresponding grouped right-hand sides.
Test each candidate against all the grouped inequalities, and keep
the first feasible candidate. This succeeds: the profile polyhedron
is nonempty and bounded by $0\le w_j\le\lambda_j$, so has a
vertex; at a vertex the active normals span $\R^m$, even if the
polyhedron is lower dimensional. If they failed to span, a sufficiently
small displacement in a nonzero orthogonal direction and its negative
would remain feasible, contradicting that it is a vertex. Thus the
enumeration includes that vertex. Its number of candidates and all
candidate tests are bounded by $2^{O(m^2)}$.

Set $w_0=t-\sum_{j=1}^mw_j$. In each gadget, the remaining
$U_i$ entries must sum to
$R_i-\sum_{B_i}w_j\in[0,\sum_{U_i}w_j]$. Fill these entries
in any fixed order, repeatedly taking the smaller of the remaining
sum and $w_j$. This takes $O((m+1)L)$ arithmetic operations and
gives the disaggregation in \cref{prop:chain-profile}. Divide
positive-weight states by their weights to obtain at most $m+1$
graph points, omitting zero states. Writing all state-arc entries
has this same $O((m+1)L)$ output scale.

All bases have integer entries of magnitude at most one. Their
inverses have cofactors bounded by $\Delta^{01}_m$ and nonzero integer
determinants bounded by $\Delta^{01}_m$. Grouping selects input-derived
affine right-hand sides; circuit tests and inverse-basis recovery
therefore use numbers with polynomial encoding length. Greedy filling
uses only additions, comparisons, and minima of these rational values.
Normalization by positive rational weights preserves polynomial
encoding length. The parameter dependence of the arithmetic count is
exponential, but the bit-length assertion needs no unit-cost assumption
for unbounded-precision numbers.
\end{proof}

\begin{remark}[The unreduced signed-subset library]\label{rem:unreduced-chain}
Keeping all $d=m+1$ profile coordinates yields normals
$\{\pm\one_S:\varnothing\ne S\subseteq\states\}$.
The same proof gives the bound $(d+1)\Delta^{01}_d$ and
$2^{O(d^2)}$ parameter work. Its positive-circuit counts for
$d=1,2,3$ are respectively $1,5,41$, with largest primitive
weights $1,1,2$. These counts are for the full signed-subset universe,
not for the rows that occur in a particular instance. Eliminating the
residual coordinate gives the smaller universe
\eqref{eq:chain-normal-universe}: for $m=1,2,3$ it has $2,6,11$
normals and $1,5,16$ positive circuits, respectively, again with
largest weights $1,1,2$. The finite enumeration in the proof reproduces
these counts exactly; the proof does not assume them. In particular,
a chain with two explicit states does not need the 41-circuit library.
\end{remark}

\subsection{Two explicit states: five tests and unit coefficients}

For $m=2$, the reduced normal universe is just
$\{\pm e_1,\pm e_2,\pm(e_1+e_2)\}$. After grouping, the
profile region takes the form
\begin{equation}\label{eq:chain-hexagon}
 \ell_1\le w_1\le u_1,\qquad
 \ell_2\le w_2\le u_2,\qquad
 \ell_s\le w_1+w_2\le u_s.
\end{equation}
All six endpoint groups are present because of the explicit and
residual weight bounds. This is the three-direction geometry of the
theta oracle, but here a single profile region describes all serial
gadgets together.

\begin{theorem}[Unit coefficients and a five-test chain oracle]\label{thm:two-state-chain}
On $G_L$ with at most two explicit simplex variables and arbitrary
observations, the sparse hull has a finite exact description with
all flow and product coefficients in $\{-1,0,1\}$.
With two explicit variables, after the original-domain and zero-row
checks, membership is equivalent to the five conditions
\begin{equation}\label{eq:chain-five-tests}
 \ell_1\le u_1,\quad \ell_2\le u_2,\quad
 \ell_s\le u_s,\quad
 \ell_1+\ell_2\le u_s,\quad
 \ell_s\le u_1+u_2.
\end{equation}
Exact separation and constructive decomposition take
$O(L+|\Obs|+1)$ rational arithmetic operations, with no circuit
enumeration and no online linear program.
\end{theorem}
\begin{proof}
The possible sum of two interval variables is
$[\ell_1+\ell_2,u_1+u_2]$. Thus \eqref{eq:chain-five-tests}
are exactly the conditions that both intervals are nonempty and
this sum interval meets $[\ell_s,u_s]$. On a feasible query set
\[
 s_* =\max(\ell_s,\ell_1+\ell_2),\qquad
 w_1=\max(\ell_1,s_*-u_2),\qquad w_2=s_*-w_1.
\]
The tests ensure $s_*\le\min(u_s,u_1+u_2)$, and these choices
satisfy every bound in \eqref{eq:chain-hexagon}. Complete the
flows by the greedy construction in \cref{thm:fixed-chain}.

Unit circuit weights alone do not exclude repeated occurrences of the
same product, so we check the original coefficients. The five
positive circuits are the three opposite pairs and
\begin{equation}\label{eq:chain-five-circuits}
 \{e_1,e_2,-\one\},\qquad
 \{-e_1,-e_2,\one\},
\end{equation}
all with weights one. The following occurrence analysis applies to
any affine row selected in each normal group. No circuit contains
both $-e_j$ and $+e_{3-j}$, or both $+e_j$ and $+\one$.

For an $a_i$ product in an only-$a_i$ state $j$, its negative
occurrences are in the local row $-e_j$ and the endpoint row
$+\one_{B_i}$; since $j\notin B_i$, the latter is either zero
or $+e_{3-j}$. Its only positive occurrence is in the other endpoint
row $+\one_{A_i\cup T_i}$. Thus a circuit cannot select two
negative occurrences. If state $j$ is doubly observed, the positive
local row $+e_j$ is also available. The other positive normal
contains $j$, so is either the same normal $+e_j$ or $+\one$;
one row is selected per normal, and these two distinct normals do
not occur together in a circuit. The negative occurrences remain
as above. Hence every $a_i$ product has coefficient of magnitude
at most one.

For a $b_i$ product in an only-$b_i$ state $j$, the negative
occurrences are $-e_j$ and $+\one_{A_i\cup T_i}$, whose latter
normal excludes $j$; the positive occurrence is the $B_i$ endpoint.
The same argument applies. A doubly observed $b_i$ product occurs
only in the two opposite local singleton rows, with opposite signs,
because $R_i$ uses the observed $a_i$ entry in that state.
A bypass product likewise occurs only in two opposite singleton rows.
Each zero-row inequality separately has unit coefficients.

The coordinate $x_{a_i}$ occurs only in its two endpoint rows,
once with each sign, and $x_{b_i}$ occurs only in the original
flow precheck. The shared coordinate $x_h$ needs an additional
check. In a positive singleton row its coefficient is zero or
$-1$, in a negative singleton row it is zero, in the positive
full row it is zero or $-1$, and in the negative full row it is
always $+1$. These statements follow by inspecting
\eqref{eq:chain-residual-reduced} and
\eqref{eq:chain-endpoints-reduced}; for example a full $B_i$
endpoint has zero $x_h$ coefficient, whereas a full $A_i\cup T_i$
endpoint has coefficient $-1$.
For the first circuit in \eqref{eq:chain-five-circuits}, two
negative contributions from its positive singleton rows are offset by the compulsory
$+1$ from the negative full row. All opposite pairs and the
second circuit also have coefficient magnitude at most one.
Consequently every affine branch of the five tests is a unit
flow/product inequality, proving the claimed finite description.

For $m=1$ the reduced system is an interval. Alternatively, embed
it as the coordinate section $y_2=0$ of the two-state instance
with no observations in label 2. Disaggregation forces state 2
to zero, so restriction of the established unit description is
exact and preserves its flow/product coefficients. The case $m=0$
is the original flow description. Grouping rows, checking the five
tests, and filling the state flows have the asserted linear cost.
\end{proof}

\subsection{Three explicit states and a sharp state-count threshold}

The unit-coefficient conclusion extends to three explicit states. Some
direct circuit sums then have bypass-flow coefficient two, but adding a
single original flow equation removes it. This is a valid change of
the ambient inequality representative; by contrast, no such change can
alter the free coefficients in the coordinate sections of
\cref{lem:section-facet}.

\begin{theorem}[Sharp unit-coefficient threshold for flat chains]\label{thm:three-state-chain}
For arbitrary observations on $G_L$, three explicit simplex variables
still admit a finite exact original-coordinate description with unit
flow and product coefficients. At most sixteen fixed positive-circuit tests,
with the original-domain and zero-row checks, give exact separation
and constructive recovery in $O(L+|\Obs|+1)$ rational arithmetic
operations. Four explicit variables already allow a hull with no
such unit-coefficient description.
\end{theorem}
\begin{proof}
For $m=3$, the positive circuits of \eqref{eq:chain-normal-universe}
have the following four forms. Each listed normal has weight one
unless specified otherwise:
\begin{align}
 &\{\one_S\}\cup\{-e_j:j\in S\}
       &&(\varnothing\ne S\subseteq[3]);\label{eq:chain3-circuits-a}\\
 &\{-\one\}\cup\{\one_S:S\in\pi\}
       &&(\pi\text{ a partition of }[3]);\label{eq:chain3-circuits-b}\\
 &\{-\one,-e_i,e_i+e_j,e_i+e_k\}
       &&(\{i,j,k\}=[3]);\label{eq:chain3-circuits-c}\\
 &\{-\one,e_1+e_2,e_1+e_3,e_2+e_3\},
       &&\text{weight two on }-\one.\label{eq:chain3-circuits-d}
\end{align}
Their numbers are $7,5,3,1$. We check completeness directly.
Without $-\one$, a positive dependence decomposes into the
dependencies of each positive subset on its negative singleton
coordinates, so minimality leaves exactly
\eqref{eq:chain3-circuits-a}. With $-\one$, presence of $+\one$
forces the opposite pair. Otherwise there are at most three other
normals, by the circuit support bound. Two negative singleton
normals would leave at most one positive proper subset, which
cannot cover all three coordinates and also supply an excess in
either negative singleton coordinate. With one negative singleton
$-e_i$, two positive subsets are necessary. The opposite singleton
$+e_i$ is excluded by minimality. Let $c>0$ be the weight on
$-\one$. The positive cover has value $c$ at $j,k$ and a strictly
larger value at $i$. Both positive subsets must contain $i$: otherwise
the only one containing $i$ would have weight greater than $c$ and,
being different from $\{i\}$, would overcover another coordinate.
Thus they are $\{i,j\}$ and $\{i,k\}$, each with weight $c$;
the weight on $-e_i$ is also $c$. This gives
\eqref{eq:chain3-circuits-c}.

With no negative singleton, at most three positive subsets form a
minimal cover with the same weighted value $c>0$ at every coordinate.
If $\{i\}$ is present, then $\{j,k\}$ either completes a partition
or is absent: its presence with any further set would contain the
proper positive dependence $\{-\one,e_i,e_j+e_k\}$.
If both pairs containing $i$ occur, their positive weights sum to
more than the weight of either pair, so the singleton makes the cover
at $i$ larger than at $j$ or $k$. If exactly one such pair, say
$\{i,j\}$, occurs, covering $k$ requires the last available set
to be $\{k\}$; the cover at $i$ then exceeds that at $j$.
Hence no pair containing $i$ occurs, leaving the singleton partition
or $\{\{i\},\{j,k\}\}$. With only pairs, two distinct pairs
overcover their common coordinate, so all three are needed. Their
coordinate equations force equal weights and weight two on $-\one$
in primitive normalization, giving \eqref{eq:chain3-circuits-d}.
This also verifies positivity and minimality of the displayed list.

All positive-normal weights are one. No circuit contains both
$-e_j$ and a positive subset normal excluding $j$, or both
$+e_j$ and a different positive subset normal containing $j$.
These facts follow immediately from the four forms. The occurrence
argument in \cref{thm:two-state-chain} therefore proves that every
product coefficient has magnitude at most one for every choice of
affine branches. The only doubled circuit weight belongs to
$-\one$, whose right-hand side $\lambda_0-t$ has no product
coordinate. Each $x_{a_i}$ has at most its two opposite-signed
endpoint occurrences, each with weight one, and no $x_{b_i}$
occurs in a profile right-hand side.

It remains to handle $x_h$. Negative singleton rows have coefficient
zero; the negative full row has coefficient $+1$. A positive
subset row has coefficient zero or $-1$. Thus circuits
\eqref{eq:chain3-circuits-a} and \eqref{eq:chain3-circuits-c}
have unit $x_h$ coefficients. In
\eqref{eq:chain3-circuits-b}, the only exception is coefficient
$-2$: the partition must consist of three singletons and all three
chosen right-hand sides must be upper endpoints $t-R_i$.
They belong to distinct gadgets, since each gadget supplies only
one upper endpoint. Choose one of them. Its $x_{a_i}$ coefficient
is $-1$, with no other contribution from that gadget. In the
circuit inequality written as an affine expression at least zero,
add the identically zero flow expression
\[
 x_{a_i}+x_{b_i}+x_h-1.
\]
The coefficients of $x_h,x_{a_i},x_{b_i}$ become $-1,0,1$.

In \eqref{eq:chain3-circuits-d}, the only exception is coefficient
$+2$, occurring when all three positive pair rows have zero
$x_h$ coefficient. A positive pair row can only come from a
gadget endpoint, so these are three lower endpoints $R_i$ from
distinct gadgets. Choose one and subtract its flow expression
$x_{a_i}+x_{b_i}+x_h-1$. The three affected coefficients become
$1,0,-1$. No other coefficient changes. These replacements preserve
validity on the original flow domain and leave the violation at any
query that passes the flow precheck unchanged. Hence every branch
of the exact circuit description has a unit-coefficient representative.

The circuit and inverse-basis libraries now have fixed size.
\Cref{thm:fixed-chain} supplies linear-cost separation and recovery,
including zero weights and lower-dimensional profile regions.
Finally \cref{ex:small-sp} has four explicit states and a local
section normal $(2,1)$. By \cref{lem:section-facet}, this ratio
cannot be removed by any affine-hull equation, so no unit-coefficient
description exists. Additional unused explicit variables can be
fixed to zero if a counterexample with more than four is desired.
\end{proof}

\begin{example}[Both bypass-coefficient repairs are needed]\label{ex:chain3-repairs}
Take $L=3$, $x_h=1/2$, and first let $y_1=y_2=y_3=1/3$,
$x_{a_i}=2/5$, $x_{b_i}=1/10$. Observe only $(a_i,i)$,
with all three products zero. The singleton-partition certificate is
\[
 \sum_{i=1}^3(t-R_i)+(\lambda_0-t)
   =2t-\sum_{i=1}^3x_{a_i}+\sum_{i=1}^3z_{a_i,i}+\lambda_0
   =-1/5.
\]
It has $x_h$ coefficient $-2$; adding any selected gadget balance
gives a unit-coefficient violated inequality.
For the other repair, take $y_1=y_2=y_3=1/4$,
$x_{a_i}=1/20$, $x_{b_i}=9/20$. Observe only $b_i$ products,
on state pairs $\{1,2\},\{1,3\},\{2,3\}$ in gadgets 1, 2, 3,
respectively, each with value $1/20$. Now $R_i=3/20$ and the
pair-triangle certificate is
\[
 R_1+R_2+R_3+2(\lambda_0-t)=-1/20.
\]
Its $x_h$ coefficient is $+2$; subtracting a selected gadget
balance gives the unit form. Both query points satisfy all original
flow, simplex, and separate McCormick inequalities. The exact
joint cuts detect the missing common profile.
\end{example}

\begin{corollary}[Only observed labels count]\label{cor:chain-observed-labels}
Let $J=\{j:(e,j)\in\Obs\text{ for some }e\}$ and $a=|J|$.
All fixed-state chain conclusions apply with $a$ in place of $m$.
In particular, at most three observed labels guarantee unit flow/product
coefficients, regardless of the number of unobserved simplex labels;
at most two observed labels permit the five-test oracle.
General separation and compact recovery take
$O((a+1)L+|\Obs|+m)+2^{O(a^2)}$ rational arithmetic operations.
\end{corollary}
\begin{proof}
Apply the exact state merger in \cref{thm:block-state-reduction},
retaining the labels in $J$ and merging all others with the residual
state. The merged weight is $1-\sum_{j\in J}y_j$.
Keep the original simplex constraints on all $m$ coordinates and
apply the preceding theorems to this reduced state list. Substituting
the merged weight changes no flow or product coefficient. The
proportional refinement \eqref{eq:proportional-refinement} recovers
each unused label from the same normalized merged flow. It can be
stored as one flow and its constituent weights, so the stated
compact-recovery cost includes only $O(m)$ additional bookkeeping.
Writing every entry of all global state flows instead can require
$\Theta((m+1)|E|)$ output operations. If $a=0$, the hull is simply
$\Flow\times\Delta_m$ and no profile elimination is needed.
\end{proof}

These theorems are specific to a chain of parallel pairs with one bypass.
A general nested series--parallel network can require several interacting
profiles, and neither the five-test result nor the fixed-state
determinant bound asserts a description for that larger class.
Conversely, \cref{thm:fibonacci} shows that no coefficient bound
independent of the state count exists even on this flat family.
